\documentclass[10pt,a4paper]{amsart}
\usepackage[margin=19mm]{geometry}
\usepackage{amsmath,amssymb,mathtools}
\usepackage[scr=rsfs,cal=euler]{mathalfa}
\usepackage{xfrac}
\usepackage[unicode,hidelinks,bookmarks=false]{hyperref}
\newcommand{\ie}{i.e.\ }
\newcommand{\cf}{cf.\ }
\newcommand{\resp}{respectively\ }
\newcommand{\Subsection}{\subsection}
\providecommand{\nobreakdash}{\nobreak\hbox{-}}
\setlength{\emergencystretch}{2em}
\multlinegap0pt
\usepackage{siunitx}
\usepackage{siunitx}
\usepackage{siunitx}
\usepackage{siunitx}
\usepackage{mathtools}
\usepackage{amssymb}
\usepackage[shortlabels]{enumitem}
\usepackage{float}
\usepackage{xcolor}
\usepackage{tikz}
\usepackage{cancel}
\usepackage{siunitx}
\newtheorem{prop}{Proposition}
\newtheorem{lemma}{Lemma}
\title{The quantum integrable hierarchy for the Gromov-Witten theory of~ elliptic~curves}
\author{Paolo Rossi, Sergey Shadrin, Ishan Jaztar Singh}
\date{Source: arXiv:2512.04621v1 (CC BY 4.0)}
\begin{document}
\maketitle

\noindent\emph{Source and licence.} The quantum integrable hierarchy for the Gromov-Witten theory of~ elliptic~curves. Version arXiv:2512.04621v1 (CC BY 4.0), distributed by arXiv under the Creative Commons Attribution 4.0 International licence, http://creativecommons.org/licenses/by/4.0/, as recorded in the arXiv licence metadata for this exact version and shown on the abstract page https://arxiv.org/abs/2512.04621v1. Full licence notice: \texttt{licenses/arxiv-2512.04621-CC-BY-4.0.txt}.\par\smallskip

\noindent\textit{Editorial scope.} Counted unit: the article's prop vanishing (source0.tex lines 445--452). The article's complete original proof of it is delivered with it (source0.tex lines 454--491, one \texttt{proof} environment of 38 lines). No mathematics is written, repaired or re-derived: the statement and the proof are the article's own lines, in the article's own notation, and no retained line is substituted or re-flowed.  No further source-local context is needed: every cross-reference of every retained line resolves inside the two delivered spans.  Nothing was omitted from any retained span, so the package carries no omission record.  No retained line contains a citation, so the wrapper carries no bibliography and no bibliography entry.  No retained line contains a figure environment, an image-inclusion command or drawing code, so no image is delivered and the package declares no asset dependency.  Editorial formatting only, in addition: an AMS \texttt{amsart} preamble that restates the article's own declarations and macros used by the retained lines, namely the environments \texttt{prop}, \texttt{lemma} on separate counters, together with the standard packages those lines need.  The retained environments print in this document as Proposition 1, Lemma 1 in order of appearance, whereas the article prints the same items with the numbering of its own full document.  The complete native source of the article as distributed in the arXiv e-print for this exact version is bundled unchanged as \texttt{sources/190338/source0.tex}.
\begin{prop}\label{prop: vanishing}
Let $(a_1,\dots,a_n)\in\mathbb{Z}^n$ with $\sum_i a_i=0$ and $g\geq 2$.
%Fix generic points $p_1,\dots,p_n\in E$ (the point class insertions $e_4=\mathrm{pt}$). 
Then
\[ 
\int_{\overline{M}_{g,n}} \mathrm{DR}_g(a_1,\ldots,a_n) \; \lambda_{g-2}\;c_{g,n}\!\big(e_4^{\otimes n}\big) =0.
\]
\end{prop}

\begin{proof} If $a_1=\cdots=a_n=0$, then  $\mathrm{DR}_g(a_1,\ldots,a_n) = (-1)^g\lambda_g$ and the desired vanishing follows from the vanishing of $\lambda_g \, c_{g,n}\big(\bigotimes_{i=1}^n e_{\alpha_i}\big)$ for $g\geq 1$ that we recall above.
	
Assume that not all $a_i$ vanish. Then we prove a more refined statement, using the following interpretation of the product $\mathrm{DR}_g(a_1,\ldots,a_n) \;c_{g,n}\!\big(e_4^{\otimes n}\big)$. As before, let $A_+$ denote the subtuple of $A$ consisting of all the positive integers in $A$, and let $A_-$ denote the subtuple of $A$ consisting of all the negative integers in $A$, and let $n_0$ be the number of $a_i$'s that are equal to $0$. Namely, let $\overline{M}^\sim_{g,n_0}(E\times \mathbb{P}^1, d, A_-, A_+)$ be the moduli space of stable relative maps of connected genus $g$ curves to the rubber trivial bundle over $E$, whose projection to $E$ has degree $d$. By relative we mean relative to the zero and infinity sections, with ramification profiles given by $A_-$ and $A_+$, respectively. Let
\begin{equation*}
	\mathsf{src} \colon \overline{M}^\sim_{g,n_0}(E
	\times \mathbb{P}^1, d, A_-, A_+) \rightarrow \overline{M}_{g,n}
\end{equation*}
be the source map that forgets the stable relative map and retains the stabilization of the source curve. Then we have
\begin{align*}
& \int_{\overline{M}_{g,n}} \mathrm{DR}_g(a_1,\ldots,a_n) \; \lambda_{g-2}\;c_{g,n}\!\big(e_4^{\otimes n}\big) 
\\ & =
\sum_{d=0}^\infty q^d \int_{\overline{M}_{g,n}} \mathsf{src}_* \left(
\left[\overline{M}^\sim_{g,n_0}(E
\times \mathbb{P}^1, d, A_-, A_+)\right]^{\mathrm{vir}} \prod_{i=1}^n \mathsf{ev}_i^*(e_4)\right) \lambda_{g-2}
,
\end{align*}
where $\mathsf{ev}_i\colon \overline{M}^\sim_{g,n_0}(E
\times \mathbb{P}^1, d, A_-, A_+) \to E$ is the evaluation map at the $i$-th marked point for the relative stable map combined with the projection to $E$. 
	
Now the statement of the proposition follows from the following lemma:
\begin{lemma} Assume that not all $a_i$ vanish. Then for any $g,d\geq 0$ we have 
\[
\left[\overline{M}^\sim_{g,n_0}(E
\times \mathbb{P}^1, d, A_-, A_+)\right]^{\mathrm{vir}} \prod_{i=1}^n \mathsf{ev}_i^*(e_4) = 0.
\]
\end{lemma}

The proof of the lemma is based on the following general observation. Let $F\colon C_1\to C_2$ be a holomorphic map of smooth compact curves and $\sum_{i=1}^n b_i p_i$ be a principal divisor on $C_1$. Let us prove that then $\sum_{i=1}^n b_i F(p_i)$ is a principal divisor on $C_2$. If $F$ is constant map, then the statment is obvious. Assume $F$ is non-constant and let $\sum_{i=1}^k b_i p_i = (f)$, where $f\colon C_1\to \mathbb{P}^1$ is a meromorphic function. Then $\sum_{i=1}^k b_i F(p_i) = (g)$, with $g\colon C_2\to \mathbb{P}^1$ defined as $g(q) \coloneqq \prod_{p\in F^{-1}(q)} f(p)^{\mathrm{ord}_p F}$, where $\mathrm{ord}_p F$ denotes the local ramification degree of $F$ at $p\in C_1$. 
	
This observation implies that for a choice of points $p_1,\dots,p_n\in E$ such that $\sum_{i=1}^n a_i p_i$ is not a principal divisor on $E$, the geometric intersection $\cap_{i=1}^n \mathsf{ev}_i^{-1}(p_i)$ on $\overline{M}^\sim_{g,n_0}(E
\times \mathbb{P}^1, d, A_-, A_+)$ is empty. Indeed, let $(C,x_,\dots,x_n,F,f)\in \overline{M}^\sim_{g,n_0}(E
\times \mathbb{P}^1, d, A_-, A_+)$ be a possibly nodal curve $C$ with marked points $x_1,\dots,x_n$ whose normalization has $l$ components $C_1,\dots,C_l$, equipped with a map $F$ to $E$ whose restriction to $C_j$ is denoted by $F_j$, $j=1,\dots,l$, and a relative stable map $f$ to the rubber $\mathbb{P}^1$, whose restriction to $C_j$ is denoted by $f_j$, $j=1,\dots,l$. Note that
\begin{itemize}
	\item The divisor of $f_j$ on $C_j$ is principal and supported on nodes and marked points. Applying the observation above to the map $F_j\colon C_j\to E$, we conclude that the image of the divisor of $f_j$ under the map $F_j$ is a principal divisor on $E$, $j=1,\dots,l$. 
	\item Each marked point $x_i$ enters the divisor of exactly one $f_j$ with coefficient $a_i$. Each node enters the divisors of exactly two $f_j$'s, with the opposite coefficients. Hence, the sum of these principal divisors is equal to $\sum_{i=1}^n a_i F(x_i) = \sum_{i=1}^n a_i \mathsf{ev}_i(C,x_,\dots,x_n,F,f)$.
\end{itemize}
Thus,  $\cap_{i=1}^n \mathsf{ev}_i^{-1}(p_i)$ is empty if  $\sum_{i=1}^n a_i p_i$ is not a principal divisor on $E$, which implies the statement of the lemma and, as a corollary, the statement of the proposition.
\end{proof}

\end{document}
